\subsection{Marked components over the ray class field and their good-prime lattices}
\label{addh:section}

Write \(k=\mathbf Q(\sqrt{-d})\), \(m=g_0g_1 15^u\) and
\(F=F_0R\) as in Proposition~\ref{zetageo:field}, with \(R\)
the ray class field of \(k\) modulo \(mO_k\). Use the chosen embedding
of \(F\) in \(\mathbf C\). The selected stages have \(d\ge7\), \(m\ge3\)
and, for the deep level, \(u\ge1\). Consider the principal marked PEL
packet of Lemma~\ref{munf:marks}, with trivial auxiliary mark and
indefinite signature \(p,q\ge1\).

\begin{proposition}[A rational CM point on each marked component]
\label{addh:constants}
Every geometric component of this packet contains a product CM point
defined over \(F\) with its prescribed mark. When \(15\mid m\), its
underlying abelian variety has good reduction at every finite place
of \(F\). Every component descends to a geometrically integral
\(F\)-component. The algebraic closure of \(F\) in its function
field, and in the function field of its complete Hodge flag bundle,
is \(F\). The same \(F\) works for all rows, with the size bound of
Proposition~\ref{zetageo:field}.
\end{proposition}
\begin{proof}
Choose an elliptic curve \(C_0/F\) with \(j=j(O_k)\) and fix a complex
identification with \(\mathbf C/O_k\). All its endomorphisms are
defined over \(F\). Indeed their actions on the tangent line give an
injective homomorphism in characteristic zero, and their scalars lie
in \(k\subset F\). As \(O_k^\times=\{\pm1\}\), its normalized Weber
function is a degree-two quotient by sign. It is defined over \(F\):
in short Weierstrass coordinates its formula is a nonzero rational
constant times \(g_2g_3 x/(g_2^3-27g_3^2)\), invariant under a
change of scale. Here \(j\ne0,1728\).

Let \(P\) correspond to \(1/m\bmod O_k\). Its Weber value is in
\(R\) by the second main theorem of complex multiplication
\cite[Theorem 11.39]{coxCM}. Every \(\sigma\in\operatorname{Gal}(\bar F/F)\)
therefore sends \(P\) to \(\epsilon(\sigma)P\) with a unique
\(\epsilon(\sigma)\in\{\pm1\}\). The signs define a continuous
character, since \(P\) has order \(m>2\). As \(P\) generates
\(C_0[m]\) over \(O_k/mO_k\), the action on the whole group is that
same sign. Twist \(C_0\) by this character using the automorphism
\([-1]\). On the resulting \(C/F\) the CM action still descends and
every \(m\)-torsion point is \(F\)-rational.

The integral \(j\)-invariant gives potential good reduction at every
finite place \(v\). When \(15\mid m\), choose \(l\in\{3,5\}\)
different from the residue characteristic at \(v\). Inertia has finite
image on \(T_lC\), and its image lies in \(1+lM_2(\mathbf Z_l)\)
by rationality of \(C[l]\). This kernel is torsion-free: if
\(U=1+l^aB\) with \(a\ge1\) and \(B\not\equiv0\bmod l\),
raising \(U\) to a positive exponent \(b\) prime to \(l\) preserves
its depth \(a\), and taking the \(l\)-th power raises that depth
by exactly one, by the binomial theorem for odd \(l\).
Inertia is therefore trivial. The
\href{https://wstein.org/papers/bib/Serre-Tate-Good_Reduction_of_Abelian_Varieties.pdf}
{N\'eron--Ogg--Shafarevich criterion, Serre--Tate, Theorem 1},
gives good reduction over \(F_v\).

For the given marked row, use the rational isometry of
Lemma~\ref{cmeig:row} and the axis point of
Proposition~\ref{cmeig:degree}. Its factors are
\(I_i\otimes_{O_k} C\), with the conjugated action on the
opposite-type factors. Each \(I_i\) is finite projective over
\(O_k\). Its idempotent on a power of \(C\), the multiplication
action of \(O_E\), and the polarized trace pairing used in that
proposition are all defined over \(F\). For the pairing one may
use the principal elliptic polarization to identify the dual
powers: every matrix entry in their homomorphisms is a CM
endomorphism of \(C\). The rigidified graph line of the polarization
is obtained by pullback of the normalized Poincar\'e bundle.
The product has constant \(m\)-torsion since it is a direct factor
of a product of curves with constant \(m\)-torsion. Thus the
transported mark is defined over \(F\), including marks mixing
the different factors. Good reduction of \(C\) also gives good
reduction of these factors, either by their integral idempotents
or by the same Tate-module criterion. For the tensor construction
see \cite[\S2]{serreTensor}.

In the marked period classification this point lies on the
prescribed geometric component. The classification of the
half-mark by the dual basis uses the root \(\zeta_m\in F\), and
the rigidified graph line and its inverse reconstruction commute
with base change. Consequently the fine moduli comparison preceding
Corollary~\ref{cmeig:final} descends to \(F\): after a finite Galois
extension each conjugate classifying map represents the same
functorial operation, so equals the original map by uniqueness.
Smooth geometric components are disjoint. The unique component
through this \(F\)-point is fixed by \(\operatorname{Gal}(\bar F/F)\)
and descends. For its function field \(K_Z\), geometric integrality
implies that \(K_Z\otimes_F\bar F\) is a domain. A nontrivial
finite extension of \(F\) inside \(K_Z\) would yield nontrivial
idempotents after this scalar extension, since \(F\) has
characteristic zero. The flag bundle has a dense affine cell over
\(K_Z\), so its function field is purely transcendental over \(K_Z\)
and has the same algebraic constants. This is the flag field in
\src{src/Formalism/LocalToGlobal/ProjectiveLattices/IsGenFrGenFr.lean}{33}{the generic flag variables}
and the condition used in
\src{src/Development/AbelianVarietyTheory/Proofs/DivisorProofs/ExistsGoodCoverOverSizedField.lean}{143}{the algebraic constants at the marked level}.
\end{proof}

\begin{proposition}[The split half-level map and its metrics]
\label{addh:level}
Over \(F\), start at a fine base level \(m_0=g_0g_1\ge3\) prime to
\(n=15^u\), \(u\ge1\), in the same classical PEL packet.
Write \(I\bar I=nO_E\) for the coprime halves in
Proposition~\ref{zetageo:field}, extended to \(O_E\), choosing \(I\)
as the annihilator of the new marked half. Forgetting the refinement
from level \(m_0n\) to \(m_0\) is finite \'etale and surjective.
On the full inverse image of the packet its degree is
\(\lvert\mathrm{GL}_r(O_E/I)\rvert\), where \(r=p+q\).
On each complex component above a base component, pullback preserves
the supremum norm and the Petersson \(L^2\)-norm formed with the
component probability measure.
\end{proposition}
\begin{proof}
On the fine base the universal abelian scheme has finite \'etale
\(n\)-torsion. The Chinese remainder idempotents split it as
\(A[I]\oplus A[\bar I]\). Both summands are isotropic and paired
perfectly with each other, since the polarization interchanges the
two idempotents under its adjoint. The first is \'etale locally a
free \(O_E/I\)-module of rank \(r\), as follows on geometric fibers
from the prescribed projective period lattice. A basis of this
summand and \(\zeta_n\) uniquely determine the dual mark. The
prime-to-\(n\) marking is retained: in the equation taking the
\(n\)-multiple of a refined torsion section to its old value,
multiplication by \(n\) is invertible on \(A[m_0]\).
Conversely every such basis extends the old level condition. At
each new local factor one lifts its matrix modulo \(I\) to an
invertible integral matrix and uses the inverse adjoint on the
conjugate factor. This supplies the compatible polarized Tate
markings at geometric points. The representing scheme is thus
the basis torsor under the indicated constant finite group.
The map on marks agrees with
\src{src/Development/AbelianVarietyTheory/BlockTheory/DivisorBase/HalfLatticeTypes.lean}{115}{the prescribed map of levels}.

The projection keeps the polarized abelian variety, tangent frame
and flag. The pointwise norm of a pulled-back section is unchanged.
On each complex component the projection is a finite covering
of its image component, and the invariant flag density of
Proposition~\ref{crev:mass} pulls back locally to the same density.
Integration multiplies by the sheet number, as does total component
volume. Division by that volume gives the claimed equality. The
supremum is unchanged by surjectivity.
\end{proof}

\begin{proposition}[Testing the subcanonical lattice at a good prime]
\label{addh:good-lattice}
Fix such a marked component, a place \(v\mid l\) of the splitting
coefficient field with \(l\nmid 2m\Delta_E\), and a smooth projective
toroidal model at its hyperspecial level there. Let \(\mathcal X\)
be its complete Hodge flag bundle over \(O_{F_v}\), \(B\) the
inverse image of the reduced toroidal boundary, and \(\mathcal L\)
the character line of an integral flag weight with central exponent
zero. Take a rational cusp rule whose evaluation on the generic
flag is a section \(s\) of \(\mathcal L(-B)_{F_v}\). Then \(s\)
belongs to \(H^0(\mathcal X,\mathcal L(-B))\) if and only if it
passes all the recorded integral-frame tests over \(v\) on this
component. It suffices to test abelian points over finite unramified
extensions of \(F_v\).
\end{proposition}
\begin{proof}
The smooth toroidal model with its universal semiabelian scheme
and relative normal crossings boundary is the good-reduction
model used in Proposition~\ref{cmeig:hyperspecial}. Its flag bundle
is smooth and proper and its interior is dense in each special
fiber component. The embedding idempotents of \(O_E\) split
integrally at \(v\). For a recorded test, extend its generic PEL
point to the toroidal model by properness. The connected
semiabelian extension agrees with the one in the test by
\href{https://grothendiecksga.com/read/sga7/en/9-3.html}
{SGA~7 I, Expos\'e IX, Proposition 3.2}. Its integral Lie
eigensummands therefore agree as well. The test bases give the
extension of the generic flag, since their successive spans are
direct summands over the test DVR. Evaluating an integral section
of \(\mathcal L(-B)\) in this frame first uses its inclusion into
\(\mathcal L\) and then an integral trivialization. This gives
an integral value, also at a degenerating test. The central
exponent zero removes the scalar theta generator from the value.

Conversely the regular model is normal, and \(s\), as a section of
\(\mathcal L(-B)\), is regular at every horizontal prime divisor
by the generic-fiber hypothesis. If \(s\) fails to extend, it has
order \(-a<0\) along a vertical prime divisor. The special fiber
is reduced, so for a uniformizer \(\pi\) the section \(\pi^a s\)
is a unit in a trivialization near its generic point. Choose a
closed point in this open set lying over the abelian interior.
Its residue field is finite. Smoothness and Hensel lifting give
a point over the unramified complete DVR with that residue field,
reducing to the chosen point. The eigenflags there admit adapted
integral bases. For its abelian scheme the principal polarization
line has a generator of its module of sections over this DVR:
on every fiber its space of sections has dimension one and higher
cohomology vanishes, so cohomology and base change apply.
These choices supply the framed abelian test. Its generic frame
is defined over a reduced finitely generated \(F\)-algebra, by
taking affine coordinates in the finite-type framed moduli space,
as required for evaluating a rule. The value of \(\pi^a s\) is a
unit, giving a test of negative valuation for \(s\). This proves
both assertions.
\end{proof}

Thus on a common rational cusp space with the indicated generic
realization, the recorded and toroidal subcanonical lattices
contribute zero local index in Lemma~\ref{zetageo:lattice} at
these primes. The test convention is the one in
\src{src/Formalism/LocalToGlobal/ProjectiveLattices/VAtVNeBot.lean}{53}{the semiabelian model and its integral eigenframes}.

In particular let \(\pi:Z'\to Z\) be one of the marked component
maps of Proposition~\ref{addh:level}. For a common rational cusp
space \(V\) on \(Z\) and its pullbacks on \(Z'\), the recorded
lattices over \(v\) satisfy
\[
       \pi^* V\ \cap\ L_{{\rm rec},v}(Z')
                    =\pi^*\bigl(V\cap L_{{\rm rec},v}(Z)\bigr).
\]
Indeed a deep test forgets to a base test of the same value.
Conversely, an abelian test detecting a negative valuation at the
base lifts to the chosen geometric component after a finite extension:
that component maps finitely and surjectively onto \(Z\). This extension
can be chosen unramified. The scheme of possible new marks over the
test DVR is the basis torsor of the direct summand \(A[I]\) of its
finite \'etale \(n\)-torsion, because \(l\nmid n\). Its generic
points split over a finite unramified extension. The lifted abelian
frame therefore detects the same negative valuation. This proves
integral descent on pullbacks even when \(l\) divides the degree
of the level map.

The comparison uses the unmodified flag line.
Exceptional divisors of the moving linear systems, the places
dividing \(2m\Delta_E\), and the final signed component height in
Corollary~\ref{recheight:period-degree} retain their own estimates.
